% Options for packages loaded elsewhere
\PassOptionsToPackage{unicode}{hyperref}
\PassOptionsToPackage{hyphens}{url}
%
\documentclass[
  11pt,
]{article}
\usepackage{amsmath,amssymb}
\usepackage{iftex}
\ifPDFTeX
  \usepackage[T1]{fontenc}
  \usepackage[utf8]{inputenc}
  \usepackage{textcomp} % provide euro and other symbols
\else % if luatex or xetex
  \usepackage{unicode-math} % this also loads fontspec
  \defaultfontfeatures{Scale=MatchLowercase}
  \defaultfontfeatures[\rmfamily]{Ligatures=TeX,Scale=1}
\fi
\usepackage{lmodern}
\ifPDFTeX\else
  % xetex/luatex font selection
\fi
% Use upquote if available, for straight quotes in verbatim environments
\IfFileExists{upquote.sty}{\usepackage{upquote}}{}
\IfFileExists{microtype.sty}{% use microtype if available
  \usepackage[]{microtype}
  \UseMicrotypeSet[protrusion]{basicmath} % disable protrusion for tt fonts
}{}
\makeatletter
\@ifundefined{KOMAClassName}{% if non-KOMA class
  \IfFileExists{parskip.sty}{%
    \usepackage{parskip}
  }{% else
    \setlength{\parindent}{0pt}
    \setlength{\parskip}{6pt plus 2pt minus 1pt}}
}{% if KOMA class
  \KOMAoptions{parskip=half}}
\makeatother
\usepackage{xcolor}
\usepackage[margin=2.2cm]{geometry}
\setlength{\emergencystretch}{3em} % prevent overfull lines
\providecommand{\tightlist}{%
  \setlength{\itemsep}{0pt}\setlength{\parskip}{0pt}}
\setcounter{secnumdepth}{-\maxdimen} % remove section numbering
\ifLuaTeX
  \usepackage{selnolig}  % disable illegal ligatures
\fi
\IfFileExists{bookmark.sty}{\usepackage{bookmark}}{\usepackage{hyperref}}
\IfFileExists{xurl.sty}{\usepackage{xurl}}{} % add URL line breaks if available
\urlstyle{same}
\hypersetup{
  pdftitle={Worksheet 5B: Right Triangle Trigonometry},
  hidelinks,
  pdfcreator={LaTeX via pandoc}}

\title{Worksheet 5B: Right Triangle Trigonometry}
\author{}
\date{}

\begin{document}
\maketitle

\begin{quote}
\textbf{What This Worksheet Is About}

Before the unit circle, trigonometry was about triangles. The ratios of
sides in a right triangle give us the original definitions of sine,
cosine, and tangent.

These ratios are the foundation for surveying, navigation, astronomy,
and engineering --- anywhere we need to find distances we can't directly
measure.
\end{quote}

\begin{center}\rule{0.5\linewidth}{0.5pt}\end{center}

\hypertarget{part-a-soh-cah-toa-the-three-ratios-20-problems}{%
\subsection{Part A: SOH-CAH-TOA --- The Three Ratios (20
problems)}\label{part-a-soh-cah-toa-the-three-ratios-20-problems}}

\hypertarget{the-definitions}{%
\subsubsection{The Definitions}\label{the-definitions}}

For a right triangle with an acute angle \(\theta\):

\[\sin\theta = \frac{\text{Opposite}}{\text{Hypotenuse}} \quad \cos\theta = \frac{\text{Adjacent}}{\text{Hypotenuse}} \quad \tan\theta = \frac{\text{Opposite}}{\text{Adjacent}}\]

\textbf{Memory aid:} \textbf{SOH-CAH-TOA} - \textbf{S}ine =
\textbf{O}pposite / \textbf{H}ypotenuse - \textbf{C}osine =
\textbf{A}djacent / \textbf{H}ypotenuse - \textbf{T}angent =
\textbf{O}pposite / \textbf{A}djacent

\begin{center}\rule{0.5\linewidth}{0.5pt}\end{center}

\textbf{For each triangle, identify the sides relative to angle
\(\theta\), then find the three ratios:}

\textbf{1.} Right triangle with legs 3 and 4, hypotenuse 5. Angle
\(\theta\) is opposite the side of length 3.

\begin{itemize}
\tightlist
\item
  Opposite = \_\_\_\_\_, Adjacent = \_\_\_\_\_, Hypotenuse = \_\_\_\_\_
\item
  \(\sin\theta ={}\) \_\_\_\_\_, \(\cos\theta ={}\) \_\_\_\_\_,
  \(\tan\theta ={}\) \_\_\_\_\_
\end{itemize}

\textbf{2.} Right triangle with legs 5 and 12, hypotenuse 13. Angle
\(\theta\) is opposite the side of length 5.

\textbf{3.} Right triangle with legs 8 and 15, hypotenuse 17. Angle
\(\theta\) is opposite the side of length 8.

\textbf{4.} Right triangle with legs 7 and 24, hypotenuse 25. Angle
\(\theta\) is opposite the side of length 24.

\begin{center}\rule{0.5\linewidth}{0.5pt}\end{center}

\textbf{Special Triangles:}

\textbf{5.} The 45-45-90 triangle has legs of equal length. If each leg
is 1: - (a) Find the hypotenuse using Pythagoras. - (b) Find
\(\sin(45°)\), \(\cos(45°)\), \(\tan(45°)\).

\textbf{6.} The 30-60-90 triangle has sides in ratio
\(1 : \sqrt{3} : 2\). - (a) If the shortest side is 1 (opposite 30°),
find the other sides. - (b) Find \(\sin(30°)\), \(\cos(30°)\),
\(\tan(30°)\). - (c) Find \(\sin(60°)\), \(\cos(60°)\), \(\tan(60°)\).

\begin{center}\rule{0.5\linewidth}{0.5pt}\end{center}

\textbf{Using a Calculator:}

\textbf{7.} Find (round to 4 decimal places): - (a) \(\sin(25°)\) - (b)
\(\cos(72°)\) - (c) \(\tan(53°)\)

\textbf{8.} Find the angle (in degrees) if: - (a) \(\sin\theta = 0.6\) -
(b) \(\cos\theta = 0.8\) - (c) \(\tan\theta = 2.5\)

\begin{center}\rule{0.5\linewidth}{0.5pt}\end{center}

\textbf{Relationships:}

\textbf{9.} In any right triangle, the two acute angles are
\textbf{complementary} (they add to 90°).

If one angle is \(\theta\), the other is \(90° - \theta\).

Show that: \(\sin\theta = \cos(90° - \theta)\)

\emph{This is why cosine is called ``co-sine'' --- it's the sine of the
complementary angle!}

\textbf{10.} Verify for \(\theta = 30°\): \(\sin(30°) = \cos(60°)\)

\begin{center}\rule{0.5\linewidth}{0.5pt}\end{center}

\hypertarget{part-b-finding-missing-sides-16-problems}{%
\subsection{Part B: Finding Missing Sides (16
problems)}\label{part-b-finding-missing-sides-16-problems}}

\hypertarget{strategy}{%
\subsubsection{Strategy}\label{strategy}}

\begin{enumerate}
\def\labelenumi{\arabic{enumi}.}
\tightlist
\item
  Label the sides relative to your angle (Opp, Adj, Hyp)
\item
  Identify what you know and what you need
\item
  Choose the ratio that uses those two sides
\item
  Set up the equation and solve
\end{enumerate}

\begin{center}\rule{0.5\linewidth}{0.5pt}\end{center}

\textbf{Find the missing side (round to 2 decimal places):}

\textbf{11.} Right triangle with angle 35°, hypotenuse 10 cm. Find the
opposite side.

\textbf{12.} Right triangle with angle 48°, adjacent side 7 m. Find the
hypotenuse.

\textbf{13.} Right triangle with angle 62°, opposite side 15 m. Find the
adjacent side.

\textbf{14.} Right triangle with angle 28°, hypotenuse 20 cm. Find the
adjacent side.

\textbf{15.} Right triangle with angle 55°, adjacent side 12 m. Find the
opposite side.

\textbf{16.} Right triangle with angle 71°, opposite side 8.5 m. Find
the hypotenuse.

\begin{center}\rule{0.5\linewidth}{0.5pt}\end{center}

** Astronomy: Measuring with Angles**

\textbf{17.} A telescope is aimed at an angle of 42° above the
horizontal to observe a star. If the telescope is 15 m from a building,
how high up the building does the line of sight hit?

\textbf{18.} The Hubble Space Telescope orbits at 547 km above Earth's
surface. From a ground station, the angle of elevation to Hubble is 28°.
- (a) Draw a diagram. - (b) What is the horizontal distance from the
station to the point directly below Hubble?

\emph{Note: This is simplified --- actual calculations account for
Earth's curvature.}

\begin{center}\rule{0.5\linewidth}{0.5pt}\end{center}

** Microscopy: Stage Tilting**

\textbf{19.} In some electron microscopes, the sample stage can be
tilted to get different views of a specimen.

If a stage holding a sample 2 mm thick is tilted at 15°, what is the
apparent thickness when viewed from directly above?

\emph{Hint: The apparent thickness is the original thickness times
\(\cos(\theta)\).}

\textbf{20.} A scanning electron microscope (SEM) detector is positioned
at 35° from the sample surface normal (perpendicular). If the detector
is 50 mm from the sample along the angled path, what is the vertical
height of the detector above the sample?

\begin{center}\rule{0.5\linewidth}{0.5pt}\end{center}

** Engineering: Ramps and Inclines**

\textbf{21.} A wheelchair ramp must have a slope no steeper than 1:12
(for every 12 units horizontal, it rises 1 unit). - (a) What angle does
this make with the horizontal? - (b) If a doorway is 0.6 m above ground
level, what is the minimum ramp length?

\textbf{22.} A road sign indicates a 6\% grade (the road rises 6 m for
every 100 m horizontal). - (a) What is the angle of incline? - (b) If a
truck travels 500 m along this road, how much altitude does it gain?

\begin{center}\rule{0.5\linewidth}{0.5pt}\end{center}

** Communications: Satellite Dishes**

\textbf{23.} A satellite dish must be aimed at a geostationary
satellite. In Dublin (latitude 53.3°N), the dish must be tilted at
approximately 26° from vertical (toward the south). - (a) What angle
does the dish make with the horizontal ground? - (b) If the dish is 1.2
m in diameter and circular, what is the vertical height of the top edge
above the bottom edge?

\begin{center}\rule{0.5\linewidth}{0.5pt}\end{center}

\hypertarget{part-c-finding-missing-angles-12-problems}{%
\subsection{Part C: Finding Missing Angles (12
problems)}\label{part-c-finding-missing-angles-12-problems}}

\hypertarget{strategy-1}{%
\subsubsection{Strategy}\label{strategy-1}}

When you know sides but need the angle, use inverse trig functions: -
\(\theta = \sin^{-1}\left(\frac{\text{Opp}}{\text{Hyp}}\right)\) or
\(\theta = \arcsin\left(\frac{\text{Opp}}{\text{Hyp}}\right)\) -
\(\theta = \cos^{-1}\left(\frac{\text{Adj}}{\text{Hyp}}\right)\) or
\(\theta = \arccos\left(\frac{\text{Adj}}{\text{Hyp}}\right)\) -
\(\theta = \tan^{-1}\left(\frac{\text{Opp}}{\text{Adj}}\right)\) or
\(\theta = \arctan\left(\frac{\text{Opp}}{\text{Adj}}\right)\)

\begin{center}\rule{0.5\linewidth}{0.5pt}\end{center}

\textbf{Find the angle \(\theta\) (round to 1 decimal place):}

\textbf{24.} Opposite = 5, Hypotenuse = 13

\textbf{25.} Adjacent = 8, Hypotenuse = 17

\textbf{26.} Opposite = 7, Adjacent = 10

\textbf{27.} Opposite = 12, Hypotenuse = 15

\textbf{28.} Adjacent = 6, Opposite = 11

\begin{center}\rule{0.5\linewidth}{0.5pt}\end{center}

** Astronomy: Parallax**

\textbf{Parallax} is the apparent shift in position of a nearby star
against distant background stars, as Earth moves from one side of its
orbit to the other.

The parallax angle \(p\) (in arcseconds) is related to the star's
distance \(d\) (in parsecs) by:

\[d = \frac{1}{p}\]

\textbf{29.} Proxima Centauri, the nearest star to our Sun, has a
parallax of 0.7687 arcseconds. - (a) Calculate its distance in parsecs.
- (b) Convert to light-years (1 parsec ≈ 3.26 light-years).

\textbf{30.} A star has a parallax of 0.01 arcseconds. - (a) What is its
distance in parsecs? - (b) Why are parallax measurements difficult for
distant stars?

\begin{center}\rule{0.5\linewidth}{0.5pt}\end{center}

** Optics: Snell's Law and Critical Angle**

When light passes from one medium to another, it bends according to
Snell's Law:

\[n_1 \sin\theta_1 = n_2 \sin\theta_2\]

When light goes from a denser to a less dense medium (like glass to
air), there's a \textbf{critical angle} beyond which light is totally
internally reflected.

\textbf{31.} Light passes from glass (\(n_1 = 1.5\)) to air
(\(n_2 = 1.0\)). - (a) If the angle in glass is 30°, find the angle in
air. - (b) The critical angle is when \(\theta_2 = 90°\). Find the
critical angle for glass-to-air.

\textbf{32.} In fiber optics, light is trapped in a glass fiber by total
internal reflection. If the glass has \(n = 1.62\): - (a) Find the
critical angle for light hitting the glass-air boundary. - (b) What does
this tell us about how the fiber must be designed?

\begin{center}\rule{0.5\linewidth}{0.5pt}\end{center}

** Crystallography: Bragg's Law**

X-ray diffraction is used to determine crystal structures. When X-rays
hit crystal planes, they reflect at angles given by Bragg's Law:

\[n\lambda = 2d\sin\theta\]

where \(\lambda\) is the X-ray wavelength, \(d\) is the spacing between
crystal planes, and \(\theta\) is the angle of incidence.

\textbf{33.} X-rays with wavelength \(\lambda = 0.154\) nm hit a crystal
with plane spacing \(d = 0.282\) nm. - (a) For first-order diffraction
(\(n = 1\)), at what angle \(\theta\) will you see a bright spot? - (b)
At what angle will you see second-order diffraction (\(n = 2\))?

\textbf{34.} In a protein crystal, you observe first-order diffraction
at \(\theta = 12.5°\) using X-rays with \(\lambda = 0.154\) nm. What is
the spacing between crystal planes?

\textbf{35.} Why can't you have \(n = 3\) diffraction if
\(\frac{\lambda}{2d} > \frac{1}{3}\)?

\emph{Hint: What's the maximum value of \(\sin\theta\)?}

\begin{center}\rule{0.5\linewidth}{0.5pt}\end{center}

\hypertarget{part-d-angles-of-elevation-and-depression-14-problems}{%
\subsection{Part D: Angles of Elevation and Depression (14
problems)}\label{part-d-angles-of-elevation-and-depression-14-problems}}

\hypertarget{definitions}{%
\subsubsection{Definitions}\label{definitions}}

\begin{itemize}
\tightlist
\item
  \textbf{Angle of elevation:} The angle UP from horizontal to an object
  above you
\item
  \textbf{Angle of depression:} The angle DOWN from horizontal to an
  object below you
\end{itemize}

\textbf{Key insight:} The angle of elevation from point A to point B
equals the angle of depression from point B to point A (alternate
interior angles with a horizontal line).

\begin{center}\rule{0.5\linewidth}{0.5pt}\end{center}

\textbf{36.} From a point 50 m from the base of a building, the angle of
elevation to the top is 65°. How tall is the building?

\textbf{37.} A 30 m tall lighthouse sits on a cliff. From a boat, the
angle of elevation to the top of the lighthouse is 12° and to the bottom
(top of cliff) is 9°. - (a) How far is the boat from the base of the
cliff? - (b) How tall is the cliff?

\textbf{38.} From the top of a 100 m cliff, the angle of depression to a
boat is 18°. How far is the boat from the base of the cliff?

\begin{center}\rule{0.5\linewidth}{0.5pt}\end{center}

** Astronomy: Measuring the Solar System**

\textbf{39.} In 1672, astronomers measured the parallax of Mars to
determine the distance to the Sun. Giovanni Cassini in Paris and Jean
Richer in Cayenne, French Guiana (6,700 km apart) observed Mars
simultaneously and found it shifted by 25.5 arcseconds against the
background stars.

If Mars was directly between Earth and the distant stars: - (a) Convert
25.5 arcseconds to degrees (1° = 3600 arcseconds). - (b) Using the
baseline of 6,700 km and this angle, estimate the distance to Mars at
that moment.

\emph{Hint: For small angles, \(\tan\theta \approx \theta\) (in
radians), so distance ≈ baseline / angle}

\textbf{40.} The Moon's orbit is inclined about 5° to Earth's orbital
plane around the Sun. This is why we don't have solar and lunar eclipses
every month.

If the Moon is 384,400 km from Earth, how far above or below Earth's
orbital plane can the Moon be at maximum?

\begin{center}\rule{0.5\linewidth}{0.5pt}\end{center}

** Microscopy: Working Distance**

The \textbf{working distance} of a microscope objective is the distance
from the front lens to the sample when in focus.

\textbf{41.} A microscope objective has a working distance of 2.1 mm and
a half-angle acceptance cone of 72°. - (a) What is the radius of the
area on the sample from which light can enter the objective? - (b) This
defines the practical ``field of view'' at the focal plane.

\textbf{42.} An electron microscope gun emits electrons in a cone with
half-angle 0.5°. If the electron source is 10 cm from the first lens: -
(a) What is the radius of the electron beam when it reaches the lens? -
(b) Why is this small angle important for electron beam quality?

\begin{center}\rule{0.5\linewidth}{0.5pt}\end{center}

** Lithography: Depth of Focus**

In lithography, the \textbf{depth of focus} (DOF) determines how much
vertical tolerance there is in placing the wafer:

\[DOF = k_2 \cdot \frac{\lambda}{NA^2}\]

This is related to trigonometry because \(NA = n\sin\theta\).

\textbf{43.} For EUV lithography with \(\lambda = 13.5\) nm and
\(NA = 0.33\), with \(k_2 = 0.5\): - (a) Calculate the depth of focus. -
(b) Why is this a challenge for manufacturing?

\textbf{44.} A hypothetical ``high-NA'' EUV system with \(NA = 0.55\): -
(a) Calculate the new depth of focus. - (b) By what factor has the DOF
decreased compared to \(NA = 0.33\)?

\begin{center}\rule{0.5\linewidth}{0.5pt}\end{center}

** Surveying: Theodolite Measurements**

\textbf{45.} A surveyor uses a theodolite (precision angle-measuring
instrument) to measure the height of a mountain. From point A, the angle
of elevation to the peak is 32°. She walks 500 m closer to the mountain
to point B, where the angle of elevation is 47°. - (a) Set up equations
using \(\tan\) for both triangles. - (b) Solve for the height of the
mountain. - (c) How far is point B from the base of the mountain?

\begin{center}\rule{0.5\linewidth}{0.5pt}\end{center}

** Navigation: Great Circle Routes**

\textbf{46.} Aircraft fly ``great circle'' routes --- the shortest path
on a sphere. The angle between two great circles at their intersection
is related to the distance and direction of travel.

Dublin (53.3°N, 6.3°W) and New York (40.7°N, 74.0°W) are connected by a
great circle route. At Dublin, the initial heading is about 252°
(roughly WSW). - (a) Why isn't the heading simply ``due west'' (270°)? -
(b) The great circle distance is about 5,120 km. Using Earth's radius
(6,371 km), what central angle does this distance subtend?

\begin{center}\rule{0.5\linewidth}{0.5pt}\end{center}

\end{document}
